\answer{ex:02:1}{長さ$3.2$~mm、間隔$25\um$の列で、検出器は$\approx129$個。発火するのは幅$v\,\tau\ln2=1.7$~mmの帯、つまり$\approx69$個で、列の半分以上。方向は帯の中心。}
\answer{ex:02:2}{窓は$-\tau\ln\frac{w_2}{\theta-w_1}\le\delta\le\tau\ln\frac{w_1}{\theta-w_2}$、すなわち$[-0.235,\,0.178]\ms$。中心は$c=\frac\tau2\ln\frac{w_1(\theta-w_1)}{w_2(\theta-w_2)}=-28$~\textmu s。方向は強い耳の側へ$28$~\textmu sずれる。分解能のほぼ3刻み分である。}
\answer{ex:02:3}{$s_iw_{ij}s_j\ge0$となる$s_i$が必要で、それはサイクルが偶数のとき、かつそのときに限り存在する。相互抑制は帰着する（持続的な振動はない）。\nt{GLUT}--\nt{GABA}ループは帰着しない（振動しうる）。}
\answer{ex:02:4}{6個のニューロン$g_k=[x+y+z\ge k]$と$\bar g_k=[\bar x+\bar y+\bar z\ge4-k]$、$k=1,2,3$。$C=g_2$、$\bar C=\bar g_2$、$S=[g_1+\bar g_2+g_3\ge2]$、$\bar S=[\bar g_1+g_2+\bar g_3\ge2]$で、計8個。ANDとORからなら$12$個。}
\answer{ex:02:5}{$T_{\min}(0.6)=0.718\,\tau$。$I_c=0.225$で、閾値付近では$T_{\min}\propto(I-I_c)^{-1/2}$。}
\answer{ex:02:6}{(a)~$500$段、ニューロン$5\cdot10^4$個。ばらつきは$4.5\ms$（$0.45\,\%$）、相対誤差は$\propto T^{-1/2}$。(b)~すべての段に共通な、ばらつき$5\,\%$のランダムな遅延係数。}
\answer{ex:02:7}{$\tau_F\ln2=0.69\,\text{s}$ごとにリフレッシュ：ニューロン1個あたり毎秒$4.3$スパイク対$20$で、$4.6$倍経済的（スパイク1個あたり$10^{-10}$~Jとして$4\cdot10^{-7}$対$2\cdot10^{-6}$~J）。フリップフロップはビットを失う。コンデンサは、リフレッシュから$0.49\,\text{s}$以内に黙らせ始めた場合にビットを保つ（確率$0.71$）。}
